\originalchapter{8.05 2012期末：对称性与粒子衰变}
原件为Fall 2012 Final Exam, 2012年12月18日13:30--16:30, 六题100分, 闭卷禁计算器. 原卷单独引用公式表, 题卷本身没有附入; 已取得课程公式表原件但不冒称逐字同一2012版. 以下为\textbf{AI独立解答}.
\originalsection{第1题：十个判断（20分）}
\begin{exercise}
判断: (1)三维中心势完整波函数必在原点为0; (2)球形无限阱有不能用角动量解释的能级简并; (3)$[H,L_z]=[H,L_x]=0$保证全旋转不变; (4)自旋半整数时轨道角动量也必须半整数才能使总角动量整数; (5)$\boldsymbol r\times\boldsymbol p=-\boldsymbol p\times\boldsymbol r$; (6)$\boldsymbol r\times\boldsymbol L=-\boldsymbol L\times\boldsymbol r$; (7)两个对易Schrödinger算符因显含时间而可使对应Heisenberg算符不对易; (8)$H_H$由$H_S$中每算符换成其Heisenberg版本得到; (9)谐振子$N_H=N_S$; (10)已知$U(t,t_0)$仍不足以求$H$.
\end{exercise}
\begin{solution}
依次\textbf{F,F,T,F,T,F,F,T,T,F}, 采用正规中心势/球阱与正常旋转生成元定义域. (1)$s$态例如氢基态原点非零, 消失的是约化径向$u(0)$. (2)球阱能量由球Bessel函数零点给, 简并$2\ell+1$由旋转解释, 不像Coulomb/振子另有系统隐藏简并. (3)Jacobi由$[L_z,L_x]=\ii\hbar L_y$推出$[H,L_y]=0$, 故与所有旋转生成元对易. (4)普通单值轨道只能整数$\ell$, 半整数自旋加整数仍半整数, 不能达到所述整数目标. (5)交叉分量交换$p_j,r_k$时$\delta_{jk}$被反对称张量消掉. (6)$\boldsymbol L\times\boldsymbol r=-\boldsymbol r\times\boldsymbol L+2\ii\hbar\boldsymbol r$. (7)相同时刻$[A_H,B_H]=U^\dagger[A_S(t),B_S(t)]U$总保持. (8)幺正共轭保持算符积及标量时间系数. (9)$[H,N]=0$. (10)$H(t)=\ii\hbar(\partial_tU)U^\dagger$, 包括整体时间相位携带的能量常数.
\end{solution}
\originalsection{第2题：指定进动与方向无关单态（10分）}
\begin{exercise}
(a)自旋$1/2$初始$+z$, 磁矩$\gamma\boldsymbol S$, 给常场使方向依次$z\to x\to y\to z$, 周期$T$. (b)将$[\ket{\boldsymbol n;+}\ket{\boldsymbol n;-}-\ket{\boldsymbol n;-}\ket{\boldsymbol n;+}]/\sqrt2$展开到$z$基, 证明任意$\boldsymbol n$都是总自旋单态.
\end{exercise}
\begin{solution}
取旋转轴$\boldsymbol n_0=(1,1,1)/\sqrt3$, 有向进动角速度$-\gamma\boldsymbol B=(2\pi/T)\boldsymbol n_0$, 即$\boldsymbol B=-2\pi\boldsymbol n_0/(\gamma T)$,$\gamma\neq0$. Rodrigues公式旋转$2\pi/3$把$z$变$x$, 再同角变$y$, 一周回$z$, 三方向都与轴等夹角. 自旋矢旋转一周虽自旋子得负号, 方向周期仍$T$.

令$c=\cos\theta/2,s=\sin\theta/2$, 方向态$\ket{\boldsymbol n;+}=c\ket++s\e^{\ii\phi}\ket-$,$\ket{\boldsymbol n;-}=-s\e^{-\ii\phi}\ket++c\ket-$. 交叉展开$++$和$--$项抵消,$+-$系数$c^2+s^2=1$,$-+$系数$-1$, 所以正好为$(\ket{+-}-\ket{-+})/\sqrt2$, 被总$S_z,S_\pm$湮灭.
\end{solution}
\originalsection{第3题：Runge--Lenz矢的垂直性（10分）}
\begin{exercise}
(a)检验$\boldsymbol r\cdot\boldsymbol L=0$, $\boldsymbol p\cdot\boldsymbol L=0$. (b)$\boldsymbol R=(\boldsymbol p\times\boldsymbol L-\ii\hbar\boldsymbol p)/(me^2)-\boldsymbol r/r$, 证明$\boldsymbol R\cdot\boldsymbol L=0$.
\end{exercise}
\begin{solution}
$r_iL_i=\epsilon_{ijk}r_ir_jp_k=0$, 因$r_ir_j$对称. $p_iL_i=\epsilon_{ijk}p_ir_jp_k$把$p_i$移过$r_j$得对称$p_ip_k$项与$\delta_{ij}$项, 都被$\epsilon$消掉. 对交叉积项,
$(\boldsymbol p\times\boldsymbol L)\cdot\boldsymbol L=\epsilon_{ijk}p_jL_kL_i=(1/2)\epsilon_{ijk}p_j[L_k,L_i]=\ii\hbar\boldsymbol p\cdot\boldsymbol L$. 因而与$-\ii\hbar\boldsymbol p$修正相消, 剩$-\boldsymbol r\cdot\boldsymbol L/r=0$. 量子修正不能先将非对易的$L_kL_i$当对称积略掉; 该修正保证$R$的正确Hermitian量子定义.
\end{solution}
\originalsection{第4题：压缩改变Hamiltonian参数（15分）}
\begin{exercise}
实$\gamma$的$S=\exp[-\gamma(a^{\dagger2}-a^2)/2]$幺正,$H(m,k)=p^2/(2m)+kx^2/2$. (a)求$S^\dagger HS=H(m',k')$. (b)比较谱和映射本征态. (c)压缩是不是原振子的对称性.
\end{exercise}
\begin{solution}
$S^\dagger xS=\e^{-\gamma}x$,$S^\dagger pS=\e^\gamma p$, 所以$m'=m\e^{-2\gamma}$,$k'=k\e^{-2\gamma}$. $\omega'=\sqrt{k'/m'}=\omega$, 两Hamiltonian幺正等谱; 原$H\ket E=E\ket E$对应新本征态$S^\dagger\ket E$. 对$\gamma\neq0$,$S^\dagger HS\neq H$, 因而不是固定$m,k$系统的守恒对称变换. 只是将它映到不同参数、相同谱的模型. 本征态的长度宽度同时改变, 能级相同并不代表算符相同.
\end{solution}
\originalsection{第5题：三维振子的隐藏算符（25分）}
\begin{exercise}
$T_{ij}=a_i^\dagger a_j$, $H=\hbar\omega(N+3/2)$. (a)写$H,N_i$. (b)求$[T_{ij},T_{pq}]$. (c)检验是否与$H$对易. (d)用一个$T$把$(a_3^\dagger)^3\ket0$变成$a_1^\dagger(a_3^\dagger)^2\ket0$. (e)在总数$N=n$组, 求最高及其他$\ell$, 用$a_L,a_R,a_3$构造最高多重态前三个$m$及次高多重态最高$m$. (f)找一个产生乘湮灭的$K$, 与$H$对易并把后者变为最高$\ell$最高$m$.
\end{exercise}
\begin{solution}
$N_i=T_{ii}$,$H=\hbar\omega(\sum_iT_{ii}+3/2)$. 用$[a_j,a_p^\dagger]=\delta_{jp}$展开并消四算符项得$[T_{ij},T_{pq}]=\delta_{jp}T_{iq}-\delta_{qi}T_{pj}$. 特别$[N,T_{pq}]=0$, 所以这些算符保持能量, 可在简并空间移动态. $T_{13}(a_3^\dagger)^3\ket0=3a_1^\dagger(a_3^\dagger)^2\ket0$, 所需为$T_{13}/3$.

本题圆基定义$a_L=(a_1+\ii a_2)/\sqrt2,a_R=(a_1-\ii a_2)/\sqrt2$, $L_+=\sqrt2\hbar(a_3^\dagger a_L-a_R^\dagger a_3)$, $L_-=L_+^\dagger$. 因而$L_z=\hbar(N_R-N_L)$, 最大$m=n$, 态$(a_R^\dagger)^n\ket0$且被$L_+$湮灭, 因而$\ell_*=n$. $N=n$的其余表示$\ell=n-2,n-4,\ldots$直到0或1. 对$\ell=n$, 未归一前三态可取
\[
\begin{aligned}
\ket{n,n}&\propto(a_R^\dagger)^n\ket0,\\
\ket{n,n-1}&\propto a_3^\dagger(a_R^\dagger)^{n-1}\ket0,\\
\ket{n,n-2}&\propto[(n-1)(a_3^\dagger)^2-a_L^\dagger a_R^\dagger](a_R^\dagger)^{n-2}\ket0.
\end{aligned}
\]
第三通式限$n\geq2$. 对$n=1$前三态另列$a_R^\dagger\ket0,a_3^\dagger\ket0,a_L^\dagger\ket0$, 对应$m=1,0,-1$, 不代入负次幂. 后两式由$L_-\!=\sqrt2\hbar(a_L^\dagger a_3-a_3^\dagger a_R)$连续作用, 整体符号可变. $n\geq2$次高最高态为$Q^\dagger(a_R^\dagger)^{n-2}\ket0$,$Q^\dagger=(a_3^\dagger)^2+2a_L^\dagger a_R^\dagger$, 它被$L_+$湮灭且$m=n-2$. 若$n=1$, 没有次高组, (f)亦无该映射对象; 前三态仍如上存在. 取$K=a_R^\dagger a_L$, 保持总数, 并有$KQ^\dagger(a_R^\dagger)^{n-2}\ket0=2(a_R^\dagger)^n\ket0$, 到归一因子即完成(f). $K$不是空间旋转角动量, 解释了不同$\ell$之间的额外简并.
\end{solution}
\originalsection{第6题：二体衰变的角动量限制（20分）}
\begin{exercise}
静止$X\to A+B$满足$\boldsymbol J=\boldsymbol S_A+\boldsymbol S_B+\boldsymbol L=\boldsymbol S_X$. (a)$s_X=0,s_A=s_B=1/2$允许哪些$\ell$. (b)$s_X=3/2,s_A=1/2,s_B=1$同问. (c)为何$n\to p+e^-$两体衰变不行（均自旋$1/2$）, 三体加反中微子却可. (d)未知$X$极化为$m_X=s_X$, 衰变$s_A=1/2,s_B=0$, 单一$\ell$, 测$P(m_A=1/2)=1/5$, $P(m_A=-1/2)=4/5$, 求$s_X,\ell$.
\end{exercise}
\begin{solution}
(a)两末态自旋先合为$S=0$或1, 要总$J=0$必须$\ell=S$, 所以$\ell=0,1$. 这只列角动量许可, 其他宇称或全同粒子限制仍可能排除某通道. (b)$S=1/2$给$\ell=1,2$,$S=3/2$给$\ell=0,1,2,3$, 合并$0,1,2,3$. (c)两个半整数自旋合成整数, 再加整数轨道仍整数, 不可能匹配初始半整数; 加另一半整数反中微子则总可为半整数, 与实际三体衰变相容.

(d)$J=\ell\pm1/2$. 上支最高态$m=\ell+1/2$只能$m_L=\ell,m_A=+1/2$, 与给定非平凡概率不符. 下支最高$m=\ell-1/2$写$a\ket{\ell,\ell-1}\ket++b\ket{\ell,\ell}\ket-$. $J_+$作用给$(a\sqrt{2\ell}+b)\hbar\ket{\ell,\ell}\ket+$, 必须0; 归一给$|a|^2=1/(2\ell+1)$,$|b|^2=2\ell/(2\ell+1)$. 因此$\ell=2,s_X=3/2$. 测概率足以定这两个数依赖题给单一$\ell$及极化最高态条件, 多通道混合时不能这样唯一反推.
\end{solution}
